\section{Mesurer la fidélité à l'image propre}
\label{sec:metrics}

L'image propre est connue dans cette expérience synthétique : elle fournit
une référence de mesure, non une information disponible dans une situation
réelle de débruitage. Plusieurs scores sont nécessaires, car une faible
erreur pixel à pixel ne garantit ni des transitions nettes ni une texture
visuellement satisfaisante. Aucune des mesures suivantes ne constitue une
mesure quantitative dédiée aux contours.

\subsection{Erreur quadratique et erreur relative}

Pour deux matrices $a$ et $b$ de même taille, comportant $P$ pixels, les
mesures implémentées sont
\begin{align}
 \operatorname{MSE}(a,b)
   &=\frac{1}{P}\sum_{i,j}(a_{i,j}-b_{i,j})^2,
   \label{eq:mse}\\
 E_{\mathrm{rel}}(a,b)
   &=\frac{\left(\sum_{i,j}(a_{i,j}-b_{i,j})^2\right)^{1/2}}
           {\left(\sum_{i,j}a_{i,j}^2\right)^{1/2}}.
   \label{eq:relative-error}
\end{align}
Dans les résultats, $a=u_{\mathrm{true}}$ et $b$ est l'observation ou un
état diffusé. Le dénominateur de \eqref{eq:relative-error} est non nul
pour les deux images étudiées. L'implémentation traite néanmoins le cas
limite : deux images nulles donnent une erreur relative nulle ; une
référence nulle et une estimation non nulle donnent $+\infty$.
Ces conventions sont testées, plutôt que masquées par une petite constante
ajoutée au dénominateur.

\subsection{PSNR et plage de référence}

Avec une plage de référence fixée à $L=1$, le rapport signal sur bruit de
crête est
\begin{equation}
 \operatorname{PSNR}(a,b)
   =10\log_{10}\!\left(\frac{L^2}{\operatorname{MSE}(a,b)}\right)
   \quad\text{en décibels}.
 \label{eq:psnr}
\end{equation}
Une MSE nulle donne $+\infty$. La plage $L$ ne dépend ni des extrema de
chaque sortie ni de ceux de l'observation bruitée : les vérités terrain
sont normalisées dans $[0,1]$ et \texttt{data\_range=1.0} est conservé
pour toutes les méthodes. Les valeurs hors de $[0,1]$ entrent donc
effectivement dans les erreurs. Cette convention rend les comparaisons
cohérentes avec la documentation officielle des métriques
\cite{skimage_metrics}.

Pour une configuration, le PSNR publié est la moyenne des dix PSNR
individuels. Il ne s'agit pas du PSNR de la MSE moyenne. Le logarithme est
non linéaire : pour des MSE positives $e_r$, la convexité de $-\log$
donne
\begin{equation}
 \frac1{10}\sum_{r=0}^{9}10\log_{10}(1/e_r)
 \ \geq\ 10\log_{10}\!\left(1/\left(\frac1{10}\sum_{r=0}^{9}e_r\right)\right).
 \label{eq:psnr-mean}
\end{equation}
L'égalité exige des erreurs identiques. L'export documentaire conserve
les colonnes agrégées du CSV sans changer cette définition.

\subsection{Le SSIM effectivement utilisé}

L'indice de similarité structurelle reprend l'idée de comparer localement
luminance, contraste et structure \cite{wang2004}. Il faut cependant
préciser sa variante numérique. Ici, chaque fenêtre valide uniforme
$7\times7$ contient $m=49$ pixels. Pour une fenêtre $W$,
\begin{align}
 \mu_a&=\frac1m\sum_{p\in W}a_p,
 &s_a^2&=\frac1{m-1}\sum_{p\in W}(a_p-\mu_a)^2,\\
 \mu_b&=\frac1m\sum_{p\in W}b_p,
 &s_{ab}&=\frac1{m-1}\sum_{p\in W}(a_p-\mu_a)(b_p-\mu_b).
\end{align}
La contribution locale et le score global sont
\begin{align}
 S_W(a,b)&=
 \frac{(2\mu_a\mu_b+C_1)(2s_{ab}+C_2)}
      {(\mu_a^2+\mu_b^2+C_1)(s_a^2+s_b^2+C_2)},
 \label{eq:ssim-local}\\
 \operatorname{SSIM}(a,b)&=\frac1{|\mathcal W|}\sum_{W\in\mathcal W}S_W(a,b),
 \qquad C_1=(0{,}01L)^2,\quad C_2=(0{,}03L)^2.
 \label{eq:ssim-global}
\end{align}
Seules les fenêtres entièrement contenues dans l'image sont moyennées :
il y en a $122\times122$ pour une image $128\times128$. Il n'y a pas de
prolongement artificiel au bord. Les covariances sont échantillonnales,
avec dénominateur $m-1$. Les moyennes locales sont calculées par sommes
cumulées NumPy. Les très petites variances négatives dues à l'annulation
en virgule flottante sont ramenées à zéro ; le score SSIM lui-même n'est
jamais borné à zéro. Une covariance négative peut produire un SSIM
négatif. Des tests indépendants vérifient notamment ce cas.

Cette définition n'est pas la configuration gaussienne $11\times11$,
d'écart-type $1{,}5$, présentée dans l'article original
\cite{wang2004}, dont les statistiques pondérées n'utilisent pas la
correction échantillonnale $m/(m-1)$. Elle est explicitement documentée et maintenue à
l'identique dans toutes les comparaisons ; la bibliothèque
\texttt{scikit-image} n'est pas utilisée pour calculer les scores de cette
exécution.

\subsection{Dispersion et limites de l'interprétation}

Pour chaque métrique $z$, les tableaux et bandes utilisent la moyenne et
l'écart-type échantillonnal des dix réalisations :
\begin{equation}
 \begin{aligned}
  \bar z&=\frac1{10}\sum_{r=0}^{9}z_r,\\
  s_z&=\left(\frac1{9}\sum_{r=0}^{9}(z_r-\bar z)^2\right)^{1/2}.
 \end{aligned}
 \label{eq:sample-aggregation}
\end{equation}
Les bandes représentent $\bar z\pm s_z$, un écart-type échantillonnal,
\emph{pas} un intervalle de confiance. Les dix réalisations renseignent
la variabilité due au bruit sur une image fixée ; elles ne mesurent pas
la variabilité entre images naturelles. Un classement numérique n'est
donc pas, à lui seul, un test de supériorité statistique ou perceptuelle.
